% Section 10.3, from lectures/L10/README.md: the objectives, the questions,
% the further reading, and what the next chapter does with all of it.


\section{Review}
\label{c10:sec:review}

\needspace{6\baselineskip}
\subsection*{What you should be able to do now}
After this chapter, you should be able to:
\begin{itemize}
  \item Name every stage of a discrete operational amplifier and say what it is responsible for.
  \item Compute a gain budget with loading, and say why one without loading is not conservative but wrong.
  \item Explain why an amplifier with no voltage gain in half its stages needs those stages.
  \item Say why an open-loop DC operating point does not exist, and what your solver does about it.
  \item Explain what a Miller capacitor is compensating and what would happen without it.
  \item Close the loop and predict the closed-loop gain, the error, and the output resistance.
  \item Run one program that predicts and solves the whole amplifier and reconciles the two.
\end{itemize}

\needspace{6\baselineskip}
\subsection*{Questions to test yourself}
You should be able to answer:
\begin{enumerate}
  \item The pair gives 1923 and the common-emitter stage gives 1923. Why is the open-loop gain not 3.7 million?
  \item The Darlington buffer has a voltage gain of 0.96. What is it worth, in decibels, and why?
  \item What is the input resistance of the output stage, and what would it be if it were a single follower instead? What does that difference cost?
  \item You solve the amplifier open-loop and the output sits at the negative rail. Is that a bug?
  \item What would happen to the closed loop if the Miller capacitor were removed?
  \item The closed-loop gain is 20 and the open-loop gain is a million. How far is the closed-loop gain from exactly 20, and what would make it closer?
  \item The output stage's input resistance depends on $h_{FE}^2$, which varies by a factor of a hundred between devices. Why does the closed-loop gain not?
\end{enumerate}

\needspace{6\baselineskip}
\subsection*{Further reading}
\begin{itemize}
  \item \emph{The Art of Electronics}, Horowitz and Hill, chapter 4x, for a walk through several real operational amplifier schematics of the kind this chapter builds.
\end{itemize}

\needspace{6\baselineskip}
\subsection*{Looking ahead}
\begin{itemize}
  \item The exam papers~(Appendix~\ref{exam:ch:about}) are four hours each and cover the whole book.
  \item What this book did not cover, and where to go for it, is in \secref{c10:sec:b8}.
\end{itemize}
